\chapter{Bienvenida}\label{cap:01-bienvenida}

\section{Qué es Lázaro Linux}

Lázaro Linux es un \termino{sistema operativo}: el programa principal de la
computadora, el que se enciende primero y deja usar todos los demás. Cumple
el mismo papel que Windows o macOS. Es gratuito, se instala desde una
memoria USB y está pensado para darle una segunda vida a computadoras y
portátiles viejos que con Windows ya van demasiado lentos o que dejaron de
recibir actualizaciones.

Con Lázaro puedes, sin instalar nada más:
\begin{itemize}
  \item navegar por internet con Firefox;
  \item escribir cartas, hacer planillas y presentaciones con LibreOffice, y
        abrir y guardar archivos de Word, Excel y PowerPoint;
  \item abrir, leer, rellenar e imprimir archivos PDF;
  \item ver YouTube sin que el video se trabe, incluso en equipos muy lentos;
  \item escuchar música, ver películas y fotos;
  \item imprimir y escanear;
  \item hacer videollamadas desde el navegador.
\end{itemize}

El nombre viene del Lázaro de la historia que vuelve a la vida. El lema,
\emph{Tu hardware aún tiene pulso}, quiere decir eso: tu equipo todavía
sirve.

\section{Linux, Alpine y Lázaro}

\termino{Linux} es una familia de sistemas operativos libres. Hay muchos,
llamados \termino{distribuciones}: Ubuntu, Linux Mint, Fedora, Debian y
cientos más. Todos comparten el mismo corazón, el \termino{kernel} Linux,
y se diferencian en los programas que traen y en cómo se ven
\cite{gnu-linux}.

Lázaro está construido sobre \textbf{Alpine Linux}, una distribución
pensada para ser pequeña, simple y segura \cite{alpine-about}. Alpine se
usa sobre todo en servidores, porque gasta muy poca memoria y muy poco
disco. Lázaro toma esa base liviana y le agrega lo que una persona
necesita en su casa: un instalador gráfico en español, un escritorio listo
para usar, los programas de todos los días y unas herramientas propias para
que el equipo funcione bien sin tocar nada.

\begin{nota}
\textbf{Libre} no solo quiere decir gratis. Quiere decir que cualquiera
puede usar el programa para lo que quiera, ver cómo está hecho, mejorarlo
y compartirlo \cite{fsf-software-libre}. Por eso Lázaro puede existir:
reúne programas libres que miles de personas hacen y mejoran en todo el
mundo.
\end{nota}

\section{Qué tiene de distinto}

\begin{description}
  \item[Hecho para equipos viejos.] Al arrancar revisa tu equipo y te
    recomienda el escritorio que mejor le va según la memoria y la tarjeta
    gráfica.
  \item[Sin pasos técnicos.] Idioma, teclado, sonido, wifi, impresoras y
    controladores de video quedan listos solos.
  \item[Funciona sin internet.] La edición Completa trae todo en el USB y
    se instala en un equipo sin red.
  \item[YouTube ligero.] Un reproductor que no usa el navegador y gasta
    mucho menos procesador (capítulo~\ref{cap:10-multimedia}).
  \item[Siempre hay un plan B.] Si el escritorio no llega a abrir, hay un
    instalador en modo texto que hace lo mismo con el teclado.
  \item[Afinado para tu equipo.] El \emph{Afinador} ajusta el sistema para
    que no se congele con muchas pestañas abiertas ni al copiar a un
    pendrive (capítulo~\ref{cap:15-herramientas}).
  \item[Avisa de las actualizaciones.] Un icono en la barra te avisa cuando
    hay novedades, y se instalan con un clic.
\end{description}

\section{¿Funciona en mi equipo?}\label{sec:requisitos}

\begin{center}
\begin{tabularx}{\textwidth}{>{\bfseries}lX}
\toprule
Procesador & De 64 bits (x86\_64). Casi todos los equipos desde 2008 lo son.
             Los muy antiguos de 32 bits (por ejemplo, muchos netbooks con
             Intel Atom N270 o N280) no sirven. \\
Memoria (RAM) & 2 GB o más es lo cómodo. Con 1 GB puede funcionar con el
             escritorio LXQt; con menos, el instalador no deja seguir. \\
Disco & 10 GB como mínimo (el instalador no deja seguir con menos); 32 GB o
             más si vas a guardar fotos y videos. \\
Arranque & BIOS o UEFI: sirven los dos, también los equipos muy viejos. \\
Memoria USB & De 4 GB para la edición Ligera o de 8 GB para la Completa. \\
\bottomrule
\end{tabularx}
\end{center}

Ya se probó en equipos reales como un Lenovo ideapad 110 (Celeron, 4 GB),
un HP Compaq CQ50 (Sempron de un núcleo) y PC con tarjetas NVIDIA de
varias épocas.

\begin{consejo}
¿No sabes cuánta memoria tiene tu equipo? En Windows, pulsa
\keys{Windows + Pausa} o busca \enquote{Acerca de} en el menú Inicio: aparece
como \enquote{RAM instalada}. Ahí mismo dice si el sistema es \enquote{de 64
bits}.
\end{consejo}

\section{Las dos ediciones}\label{sec:ediciones}

Lázaro se descarga en dos versiones, llamadas \emph{ediciones}. Las dos
instalan el mismo sistema; cambia lo que traen dentro del USB.

\begin{center}
\begin{tabularx}{\textwidth}{lXX}
\toprule
 & \textbf{Ligera} & \textbf{Completa} \\
\midrule
Tamaño & menos de 2 GB & unos 4,6 GB \\
¿Necesita internet al instalar? & Para GNOME, KDE Plasma y las aplicaciones
  de Flathub, sí & No (solo para Google Chrome) \\
Escritorios sin internet & XFCE, MATE, LXQt, Sway, XFCE Wayland & Los 7, en
  todos los idiomas \\
Controladores (firmware) & Los de red & Todos \\
\bottomrule
\end{tabularx}
\end{center}

\textbf{Si dudas, elige la Ligera}: se descarga más rápido y alcanza para la
mayoría. Elige la Completa si el equipo no va a tener internet o si tu
conexión es lenta.

\section{Antes de seguir}

Lázaro está todavía en desarrollo. Ya se instala bien en muchos equipos
reales, pero:
\begin{cuidado}
Instalar Lázaro \textbf{borra el disco que elijas}. Antes de instalar,
guarda una copia de tus fotos, documentos y todo lo que te importe en un
pendrive, un disco externo o en la nube. El capítulo~\ref{cap:02-antes-de-instalar}
explica cómo.
\end{cuidado}

Si es tu primera vez, sigue el orden del libro: el capítulo~2 prepara todo,
el~3 te deja probar Lázaro sin tocar nada del equipo, y el~4 lo instala.

\fuentes
